\begin{reportSheet}
\section{Ход работы}

\subsection{Описание модели требований}

На рисунке~\ref{fig:use-cases} приведена уточнённая диаграмма прецедентов.
Пользователь выполняет поиск, администратор запускает анализ, управляет
заданиями, результатами и параметрами моделей. Catalog Service инициирует
обработку и получает её статус, Auth Service обеспечивает аутентификацию.

\analysisDiagram[0.72]{05}{Диаграмма прецедентов}{fig:use-cases}
\end{reportSheet}

\begin{reportSheet}
\subsection{Трассировка прецедентов в кооперации}

Для 18 прецедентов определены одноимённые кооперации
(рисунок~\ref{fig:traces}). Связь \texttt{trace} направлена от кооперации
к реализуемому прецеденту. Детально рассмотрены анализ аудиозаписи,
поиск аудиозаписи и настройка параметров анализа.

\rotatedAnalysisDiagram{06}{Трассировка прецедентов в кооперации}{fig:traces}
\end{reportSheet}

\subsection{Обзорные диаграммы классов анализа}

Выделены 11 классов сущностей (рисунок~\ref{fig:entities}).
Аудиозапись связана с файлом и заданиями анализа. Задание содержит
статус, даты, прогресс, источник запроса и снимок порогов при запуске,
ему принадлежат оценки уверенности. Модель связана с профилем весов
и текущими порогами.

\analysisDiagram[0.68]{09}{Обзорная диаграмма классов сущностей}{fig:entities}
\clearpage

Граничные классы представлены восемью формами интерфейса и тремя
интерфейсами внешних систем (рисунок~\ref{fig:boundaries}). Они принимают
действия акторов и передают запросы управляющим классам.

Семь управляющих классов координируют вход, поиск, анализ, повторную
обработку каталога, работу с заданиями, результатами и параметрами
(рисунок~\ref{fig:controls}). Связанные операции нескольких прецедентов
сосредоточены в общих управляющих классах.

\fullPageAnalysisDiagram{07}{Обзорная диаграмма граничных классов}{fig:boundaries}
\fullPageAnalysisDiagram{08}{Обзорная диаграмма управляющих классов}{fig:controls}
\clearpage

\subsection{Диаграммы классов-участников коопераций}

Кооперация «Проанализировать аудиозапись» объединяет форму запуска,
контроллер анализа, интерфейс файлового хранилища, контроллер задания,
аудиозапись и задание анализа (рисунок~\ref{fig:analysis-participants}).
Контроллер анализа получает файл и определяет характеристики,
контроллер задания создаёт задание и сохраняет результат.

\rotatedAnalysisDiagram{12}{Классы-участники кооперации анализа аудиозаписи}{fig:analysis-participants}
\clearpage

В кооперации «Выполнить поиск аудиозаписи» участвуют форма поиска,
управляющий класс поиска и сущность «Аудиозапись»
(рисунок~\ref{fig:search-participants}). Форма содержит описание,
фильтры и результаты, управляющий класс формирует условия поиска.

\rotatedAnalysisDiagram{11}{Классы-участники кооперации поиска аудиозаписи}{fig:search-participants}
\clearpage

В кооперации «Настроить параметры анализа» форма передаёт выбор модели,
импорт файла и изменение порогов управляющему классу. Он обращается
к модели, профилю весов и порогам уверенности
(рисунок~\ref{fig:parameters-participants}).

\analysisDiagram[0.65]{13}{Классы-участники кооперации настройки параметров анализа}{fig:parameters-participants}

\begin{reportSheet}
\subsection{Диаграммы последовательностей}

Основной поток анализа начинается действием администратора
(рисунок~\ref{fig:analysis-sequence}). Создаётся задание, проверяется
аудиозапись и получается файл. После определения характеристик оценки
отбираются по порогам уверенности. Принятые оценки обновляют данные
аудиозаписи, результат сохраняется, задание завершается.

Поиск начинается с ввода описания и/или характеристик пользователем
(рисунок~\ref{fig:search-sequence}). Форма передаёт критерии
управляющему классу -- он формирует условия, получает подходящие аудиозаписи
и возвращает результаты форме.

Настройка включает три необязательных действия администратора:
выбор модели, импорт модели с профилем весов и изменение порогов
(рисунок~\ref{fig:parameters-sequence}). Циклы отражают создание
и обновление порогов. Сохранённые параметры применяются при последующих
запусках анализа.

\rotatedAnalysisDiagram{15}{Последовательность анализа аудиозаписи}{fig:analysis-sequence}
\fullPageAnalysisDiagram{14}{Последовательность поиска аудиозаписи}{fig:search-sequence}
\fullPageAnalysisDiagram{16}{Последовательность настройки параметров анализа}{fig:parameters-sequence}
\end{reportSheet}

\begin{reportSheet}
\subsection{Анализ классов и обзорная диаграмма}

По сообщениям последовательностей уточнены операции классов:
у форм --- обработчики действий, у управляющих классов --- шаги сценариев,
у сущностей --- операции чтения и изменения данных.
Обзорная диаграмма объединяет 29 классов: 11 граничных, 7 управляющих
и 11 сущностей (рисунок~\ref{fig:all-classes}).

\rotatedAnalysisDiagram{10}{Обзорная диаграмма всех классов анализа}{fig:all-classes}
\end{reportSheet}

\begin{reportSheet}
\subsection{Диаграммы пакетов}

Классы распределены по восьми пакетам: анализ аудиозаписей, задания
и результаты, модели и параметры, аудиозаписи, поиск аудиозаписей,
интеграция с Auth Service, интеграция с Catalog Service и доступ
к файловому хранилищу (рисунок~\ref{fig:class-packages}).

Зависимости пакетов приведены на рисунке~\ref{fig:package-dependencies}.
Анализ использует задания и результаты, модели и параметры, аудиозаписи
и файловое хранилище. Поиск зависит от аудиозаписей. Входной интерфейс
Catalog Service обращается к анализу и заданиям. Зависимости направлены
от специализированного прикладного уровня к общим и базовым пакетам.

\rotatedAnalysisDiagram{19}{Распределение классов анализа по пакетам}{fig:class-packages}
\rotatedAnalysisDiagram{20}{Зависимости пакетов анализа}{fig:package-dependencies}
\end{reportSheet}
